\section{Introduction}\label{sec:intro}

Public blockchains record every payment on an immutable ledger, yet the
parties behind the payments are pseudonymous. This combination makes
cryptocurrencies attractive for money laundering, ransomware, darknet markets
and scams; Chainalysis estimates that illicit addresses received
USD~154 billion in 2025~\cite{chainalysis2026}. Exchanges and regulators
therefore need automated transaction monitoring, and the Elliptic Bitcoin
dataset~\cite{weber2019} and its extension Elliptic++~\cite{elmougy2023} have
become the standard benchmarks for learning such monitors.

Because Bitcoin transactions form a graph of fund flows, the benchmark has been
a showcase for graph neural networks (GNNs). The evidence that GNNs actually
help, however, has steadily eroded. Weber~et~al.\ already found that a Random
Forest beat a graph convolutional network~\cite{weber2019}; recent
re-evaluations show that once test-period structure is withheld from training,
a Random Forest outperforms every GNN and randomly rewired edges work as well as
the real ones~\cite{maganti2026}, and a multi-graph model fusing the
transaction and wallet graphs still loses to gradient-boosted trees on the
transaction task~\cite{ichull2026twin}. At the same time, an audit of the
dataset itself revealed that 72 of its 93 ``local'' features are statistics of
the participating addresses computed \emph{after} the end of the test
period~\cite{safar2026}, so most published numbers---including the strong
tree baselines---are affected by look-ahead leakage.

This leaves the motivating question of graph learning on this problem
unanswered: \emph{if the transaction graph does not help, is it because graph
structure carries no signal, or because we have been giving models the wrong
graph?} We argue the latter and test it. Elliptic transactions are linked only
to other transactions of the same two-week time step, and the resulting graph
is sparse (mean degree~2.3) and only moderately homophilous. Investigators,
however, do not reason over transaction-to-transaction edges alone: they link
transactions through the \emph{addresses} and \emph{entities} that control
them~\cite{meiklejohn2013}. A knowledge graph (KG) that represents
transactions and addresses as typed entities with typed relations makes this
entity structure explicit.

\textbf{Contributions.}
\begin{enumerate}
\item We build a \emph{time-sliced transaction--address knowledge graph}
(TS-KG) from Elliptic++, in which addresses are instantiated per time step and
carry only causally computed history attributes. All edges stay inside one
time step, so the evaluation is strictly inductive by construction, and the
graph is free of the look-ahead address features identified in~\cite{safar2026}.
\item We provide a measured motivation for graph learning: the TS-KG
relations are far more homophilous than the transaction graph (class-adjusted
homophily 0.92 vs.\ 0.38 on the test period) and reach 93\% of illicit
transactions, and we show which part of this signal is an artefact of how
labels were assigned.
\item We run a leakage-controlled comparison of tree ensembles and GNNs under
one protocol, separating the effect of \emph{feature leakage}, of \emph{KG
knowledge} supplied as features, and of \emph{message passing} over the KG,
with relation-level ablations, an unseen-entity test view, imbalance
strategies, explanations, and compute cost, over multiple seeds.
\end{enumerate}
All code is released for reproduction.\footnote{\url{https://github.com/TahaNaqv/kg-bitcoin-aml}}
